\documentclass[11pt]{article}

    \usepackage[breakable]{tcolorbox}
    \usepackage{parskip} % Stop auto-indenting (to mimic markdown behaviour)
    

    % Basic figure setup, for now with no caption control since it's done
    % automatically by Pandoc (which extracts ![](path) syntax from Markdown).
    \usepackage{graphicx}
    % Keep aspect ratio if custom image width or height is specified
    \setkeys{Gin}{keepaspectratio}
    % Maintain compatibility with old templates. Remove in nbconvert 6.0
    \let\Oldincludegraphics\includegraphics
    % Ensure that by default, figures have no caption (until we provide a
    % proper Figure object with a Caption API and a way to capture that
    % in the conversion process - todo).
    \usepackage{caption}
    \DeclareCaptionFormat{nocaption}{}
    \captionsetup{format=nocaption,aboveskip=0pt,belowskip=0pt}

    \usepackage{float}
    \floatplacement{figure}{H} % forces figures to be placed at the correct location
    \usepackage{xcolor} % Allow colors to be defined
    \usepackage{enumerate} % Needed for markdown enumerations to work
    \usepackage{geometry} % Used to adjust the document margins
    \usepackage{amsmath} % Equations
    \usepackage{amssymb} % Equations
    \usepackage{textcomp} % defines textquotesingle
    % Hack from http://tex.stackexchange.com/a/47451/13684:
    \AtBeginDocument{%
        \def\PYZsq{\textquotesingle}% Upright quotes in Pygmentized code
    }
    \usepackage{upquote} % Upright quotes for verbatim code
    \usepackage{eurosym} % defines \euro

    \usepackage{iftex}
    \ifPDFTeX
        \usepackage[T1]{fontenc}
        \IfFileExists{alphabeta.sty}{
              \usepackage{alphabeta}
          }{
              \usepackage[mathletters]{ucs}
              \usepackage[utf8x]{inputenc}
          }
    \else
        \usepackage{fontspec}
        \usepackage{unicode-math}
    \fi

    \usepackage{fancyvrb} % verbatim replacement that allows latex
    \usepackage{grffile} % extends the file name processing of package graphics
                         % to support a larger range
    \makeatletter % fix for old versions of grffile with XeLaTeX
    \@ifpackagelater{grffile}{2019/11/01}
    {
      % Do nothing on new versions
    }
    {
      \def\Gread@@xetex#1{%
        \IfFileExists{"\Gin@base".bb}%
        {\Gread@eps{\Gin@base.bb}}%
        {\Gread@@xetex@aux#1}%
      }
    }
    \makeatother
    \usepackage[Export]{adjustbox} % Used to constrain images to a maximum size
    \adjustboxset{max size={0.9\linewidth}{0.9\paperheight}}

    % The hyperref package gives us a pdf with properly built
    % internal navigation ('pdf bookmarks' for the table of contents,
    % internal cross-reference links, web links for URLs, etc.)
    \usepackage{hyperref}
    % The default LaTeX title has an obnoxious amount of whitespace. By default,
    % titling removes some of it. It also provides customization options.
    \usepackage{titling}
    \usepackage{longtable} % longtable support required by pandoc >1.10
    \usepackage{booktabs}  % table support for pandoc > 1.12.2
    \usepackage{array}     % table support for pandoc >= 2.11.3
    \usepackage{calc}      % table minipage width calculation for pandoc >= 2.11.1
    \usepackage[inline]{enumitem} % IRkernel/repr support (it uses the enumerate* environment)
    \usepackage[normalem]{ulem} % ulem is needed to support strikethroughs (\sout)
                                % normalem makes italics be italics, not underlines
    \usepackage{soul}      % strikethrough (\st) support for pandoc >= 3.0.0
    \usepackage{mathrsfs}
    

    
    % Colors for the hyperref package
    \definecolor{urlcolor}{rgb}{0,.145,.698}
    \definecolor{linkcolor}{rgb}{.71,0.21,0.01}
    \definecolor{citecolor}{rgb}{.12,.54,.11}

    % ANSI colors
    \definecolor{ansi-black}{HTML}{3E424D}
    \definecolor{ansi-black-intense}{HTML}{282C36}
    \definecolor{ansi-red}{HTML}{E75C58}
    \definecolor{ansi-red-intense}{HTML}{B22B31}
    \definecolor{ansi-green}{HTML}{00A250}
    \definecolor{ansi-green-intense}{HTML}{007427}
    \definecolor{ansi-yellow}{HTML}{DDB62B}
    \definecolor{ansi-yellow-intense}{HTML}{B27D12}
    \definecolor{ansi-blue}{HTML}{208FFB}
    \definecolor{ansi-blue-intense}{HTML}{0065CA}
    \definecolor{ansi-magenta}{HTML}{D160C4}
    \definecolor{ansi-magenta-intense}{HTML}{A03196}
    \definecolor{ansi-cyan}{HTML}{60C6C8}
    \definecolor{ansi-cyan-intense}{HTML}{258F8F}
    \definecolor{ansi-white}{HTML}{C5C1B4}
    \definecolor{ansi-white-intense}{HTML}{A1A6B2}
    \definecolor{ansi-default-inverse-fg}{HTML}{FFFFFF}
    \definecolor{ansi-default-inverse-bg}{HTML}{000000}

    % common color for the border for error outputs.
    \definecolor{outerrorbackground}{HTML}{FFDFDF}

    % commands and environments needed by pandoc snippets
    % extracted from the output of `pandoc -s`
    \providecommand{\tightlist}{%
      \setlength{\itemsep}{0pt}\setlength{\parskip}{0pt}}
    \DefineVerbatimEnvironment{Highlighting}{Verbatim}{commandchars=\\\{\}}
    % Add ',fontsize=\small' for more characters per line
    \newenvironment{Shaded}{}{}
    \newcommand{\KeywordTok}[1]{\textcolor[rgb]{0.00,0.44,0.13}{\textbf{{#1}}}}
    \newcommand{\DataTypeTok}[1]{\textcolor[rgb]{0.56,0.13,0.00}{{#1}}}
    \newcommand{\DecValTok}[1]{\textcolor[rgb]{0.25,0.63,0.44}{{#1}}}
    \newcommand{\BaseNTok}[1]{\textcolor[rgb]{0.25,0.63,0.44}{{#1}}}
    \newcommand{\FloatTok}[1]{\textcolor[rgb]{0.25,0.63,0.44}{{#1}}}
    \newcommand{\CharTok}[1]{\textcolor[rgb]{0.25,0.44,0.63}{{#1}}}
    \newcommand{\StringTok}[1]{\textcolor[rgb]{0.25,0.44,0.63}{{#1}}}
    \newcommand{\CommentTok}[1]{\textcolor[rgb]{0.38,0.63,0.69}{\textit{{#1}}}}
    \newcommand{\OtherTok}[1]{\textcolor[rgb]{0.00,0.44,0.13}{{#1}}}
    \newcommand{\AlertTok}[1]{\textcolor[rgb]{1.00,0.00,0.00}{\textbf{{#1}}}}
    \newcommand{\FunctionTok}[1]{\textcolor[rgb]{0.02,0.16,0.49}{{#1}}}
    \newcommand{\RegionMarkerTok}[1]{{#1}}
    \newcommand{\ErrorTok}[1]{\textcolor[rgb]{1.00,0.00,0.00}{\textbf{{#1}}}}
    \newcommand{\NormalTok}[1]{{#1}}

    % Additional commands for more recent versions of Pandoc
    \newcommand{\ConstantTok}[1]{\textcolor[rgb]{0.53,0.00,0.00}{{#1}}}
    \newcommand{\SpecialCharTok}[1]{\textcolor[rgb]{0.25,0.44,0.63}{{#1}}}
    \newcommand{\VerbatimStringTok}[1]{\textcolor[rgb]{0.25,0.44,0.63}{{#1}}}
    \newcommand{\SpecialStringTok}[1]{\textcolor[rgb]{0.73,0.40,0.53}{{#1}}}
    \newcommand{\ImportTok}[1]{{#1}}
    \newcommand{\DocumentationTok}[1]{\textcolor[rgb]{0.73,0.13,0.13}{\textit{{#1}}}}
    \newcommand{\AnnotationTok}[1]{\textcolor[rgb]{0.38,0.63,0.69}{\textbf{\textit{{#1}}}}}
    \newcommand{\CommentVarTok}[1]{\textcolor[rgb]{0.38,0.63,0.69}{\textbf{\textit{{#1}}}}}
    \newcommand{\VariableTok}[1]{\textcolor[rgb]{0.10,0.09,0.49}{{#1}}}
    \newcommand{\ControlFlowTok}[1]{\textcolor[rgb]{0.00,0.44,0.13}{\textbf{{#1}}}}
    \newcommand{\OperatorTok}[1]{\textcolor[rgb]{0.40,0.40,0.40}{{#1}}}
    \newcommand{\BuiltInTok}[1]{{#1}}
    \newcommand{\ExtensionTok}[1]{{#1}}
    \newcommand{\PreprocessorTok}[1]{\textcolor[rgb]{0.74,0.48,0.00}{{#1}}}
    \newcommand{\AttributeTok}[1]{\textcolor[rgb]{0.49,0.56,0.16}{{#1}}}
    \newcommand{\InformationTok}[1]{\textcolor[rgb]{0.38,0.63,0.69}{\textbf{\textit{{#1}}}}}
    \newcommand{\WarningTok}[1]{\textcolor[rgb]{0.38,0.63,0.69}{\textbf{\textit{{#1}}}}}
    \makeatletter
    \newsavebox\pandoc@box
    \newcommand*\pandocbounded[1]{%
      \sbox\pandoc@box{#1}%
      % scaling factors for width and height
      \Gscale@div\@tempa\textheight{\dimexpr\ht\pandoc@box+\dp\pandoc@box\relax}%
      \Gscale@div\@tempb\linewidth{\wd\pandoc@box}%
      % select the smaller of both
      \ifdim\@tempb\p@<\@tempa\p@
        \let\@tempa\@tempb
      \fi
      % scaling accordingly (\@tempa < 1)
      \ifdim\@tempa\p@<\p@
        \scalebox{\@tempa}{\usebox\pandoc@box}%
      % scaling not needed, use as it is
      \else
        \usebox{\pandoc@box}%
      \fi
    }
    \makeatother

    % Define a nice break command that doesn't care if a line doesn't already
    % exist.
    \def\br{\hspace*{\fill} \\* }
    % Math Jax compatibility definitions
    \def\gt{>}
    \def\lt{<}
    \let\Oldtex\TeX
    \let\Oldlatex\LaTeX
    \renewcommand{\TeX}{\textrm{\Oldtex}}
    \renewcommand{\LaTeX}{\textrm{\Oldlatex}}
    % Document parameters
    % Document title
    \title{geracao}
    
    
    
    
    
    
    
% Pygments definitions
\makeatletter
\def\PY@reset{\let\PY@it=\relax \let\PY@bf=\relax%
    \let\PY@ul=\relax \let\PY@tc=\relax%
    \let\PY@bc=\relax \let\PY@ff=\relax}
\def\PY@tok#1{\csname PY@tok@#1\endcsname}
\def\PY@toks#1+{\ifx\relax#1\empty\else%
    \PY@tok{#1}\expandafter\PY@toks\fi}
\def\PY@do#1{\PY@bc{\PY@tc{\PY@ul{%
    \PY@it{\PY@bf{\PY@ff{#1}}}}}}}
\def\PY#1#2{\PY@reset\PY@toks#1+\relax+\PY@do{#2}}

\@namedef{PY@tok@w}{\def\PY@tc##1{\textcolor[rgb]{0.73,0.73,0.73}{##1}}}
\@namedef{PY@tok@c}{\let\PY@it=\textit\def\PY@tc##1{\textcolor[rgb]{0.24,0.48,0.48}{##1}}}
\@namedef{PY@tok@cp}{\def\PY@tc##1{\textcolor[rgb]{0.61,0.40,0.00}{##1}}}
\@namedef{PY@tok@k}{\let\PY@bf=\textbf\def\PY@tc##1{\textcolor[rgb]{0.00,0.50,0.00}{##1}}}
\@namedef{PY@tok@kp}{\def\PY@tc##1{\textcolor[rgb]{0.00,0.50,0.00}{##1}}}
\@namedef{PY@tok@kt}{\def\PY@tc##1{\textcolor[rgb]{0.69,0.00,0.25}{##1}}}
\@namedef{PY@tok@o}{\def\PY@tc##1{\textcolor[rgb]{0.40,0.40,0.40}{##1}}}
\@namedef{PY@tok@ow}{\let\PY@bf=\textbf\def\PY@tc##1{\textcolor[rgb]{0.67,0.13,1.00}{##1}}}
\@namedef{PY@tok@nb}{\def\PY@tc##1{\textcolor[rgb]{0.00,0.50,0.00}{##1}}}
\@namedef{PY@tok@nf}{\def\PY@tc##1{\textcolor[rgb]{0.00,0.00,1.00}{##1}}}
\@namedef{PY@tok@nc}{\let\PY@bf=\textbf\def\PY@tc##1{\textcolor[rgb]{0.00,0.00,1.00}{##1}}}
\@namedef{PY@tok@nn}{\let\PY@bf=\textbf\def\PY@tc##1{\textcolor[rgb]{0.00,0.00,1.00}{##1}}}
\@namedef{PY@tok@ne}{\let\PY@bf=\textbf\def\PY@tc##1{\textcolor[rgb]{0.80,0.25,0.22}{##1}}}
\@namedef{PY@tok@nv}{\def\PY@tc##1{\textcolor[rgb]{0.10,0.09,0.49}{##1}}}
\@namedef{PY@tok@no}{\def\PY@tc##1{\textcolor[rgb]{0.53,0.00,0.00}{##1}}}
\@namedef{PY@tok@nl}{\def\PY@tc##1{\textcolor[rgb]{0.46,0.46,0.00}{##1}}}
\@namedef{PY@tok@ni}{\let\PY@bf=\textbf\def\PY@tc##1{\textcolor[rgb]{0.44,0.44,0.44}{##1}}}
\@namedef{PY@tok@na}{\def\PY@tc##1{\textcolor[rgb]{0.41,0.47,0.13}{##1}}}
\@namedef{PY@tok@nt}{\let\PY@bf=\textbf\def\PY@tc##1{\textcolor[rgb]{0.00,0.50,0.00}{##1}}}
\@namedef{PY@tok@nd}{\def\PY@tc##1{\textcolor[rgb]{0.67,0.13,1.00}{##1}}}
\@namedef{PY@tok@s}{\def\PY@tc##1{\textcolor[rgb]{0.73,0.13,0.13}{##1}}}
\@namedef{PY@tok@sd}{\let\PY@it=\textit\def\PY@tc##1{\textcolor[rgb]{0.73,0.13,0.13}{##1}}}
\@namedef{PY@tok@si}{\let\PY@bf=\textbf\def\PY@tc##1{\textcolor[rgb]{0.64,0.35,0.47}{##1}}}
\@namedef{PY@tok@se}{\let\PY@bf=\textbf\def\PY@tc##1{\textcolor[rgb]{0.67,0.36,0.12}{##1}}}
\@namedef{PY@tok@sr}{\def\PY@tc##1{\textcolor[rgb]{0.64,0.35,0.47}{##1}}}
\@namedef{PY@tok@ss}{\def\PY@tc##1{\textcolor[rgb]{0.10,0.09,0.49}{##1}}}
\@namedef{PY@tok@sx}{\def\PY@tc##1{\textcolor[rgb]{0.00,0.50,0.00}{##1}}}
\@namedef{PY@tok@m}{\def\PY@tc##1{\textcolor[rgb]{0.40,0.40,0.40}{##1}}}
\@namedef{PY@tok@gh}{\let\PY@bf=\textbf\def\PY@tc##1{\textcolor[rgb]{0.00,0.00,0.50}{##1}}}
\@namedef{PY@tok@gu}{\let\PY@bf=\textbf\def\PY@tc##1{\textcolor[rgb]{0.50,0.00,0.50}{##1}}}
\@namedef{PY@tok@gd}{\def\PY@tc##1{\textcolor[rgb]{0.63,0.00,0.00}{##1}}}
\@namedef{PY@tok@gi}{\def\PY@tc##1{\textcolor[rgb]{0.00,0.52,0.00}{##1}}}
\@namedef{PY@tok@gr}{\def\PY@tc##1{\textcolor[rgb]{0.89,0.00,0.00}{##1}}}
\@namedef{PY@tok@ge}{\let\PY@it=\textit}
\@namedef{PY@tok@gs}{\let\PY@bf=\textbf}
\@namedef{PY@tok@ges}{\let\PY@bf=\textbf\let\PY@it=\textit}
\@namedef{PY@tok@gp}{\let\PY@bf=\textbf\def\PY@tc##1{\textcolor[rgb]{0.00,0.00,0.50}{##1}}}
\@namedef{PY@tok@go}{\def\PY@tc##1{\textcolor[rgb]{0.44,0.44,0.44}{##1}}}
\@namedef{PY@tok@gt}{\def\PY@tc##1{\textcolor[rgb]{0.00,0.27,0.87}{##1}}}
\@namedef{PY@tok@err}{\def\PY@bc##1{{\setlength{\fboxsep}{\string -\fboxrule}\fcolorbox[rgb]{1.00,0.00,0.00}{1,1,1}{\strut ##1}}}}
\@namedef{PY@tok@kc}{\let\PY@bf=\textbf\def\PY@tc##1{\textcolor[rgb]{0.00,0.50,0.00}{##1}}}
\@namedef{PY@tok@kd}{\let\PY@bf=\textbf\def\PY@tc##1{\textcolor[rgb]{0.00,0.50,0.00}{##1}}}
\@namedef{PY@tok@kn}{\let\PY@bf=\textbf\def\PY@tc##1{\textcolor[rgb]{0.00,0.50,0.00}{##1}}}
\@namedef{PY@tok@kr}{\let\PY@bf=\textbf\def\PY@tc##1{\textcolor[rgb]{0.00,0.50,0.00}{##1}}}
\@namedef{PY@tok@bp}{\def\PY@tc##1{\textcolor[rgb]{0.00,0.50,0.00}{##1}}}
\@namedef{PY@tok@fm}{\def\PY@tc##1{\textcolor[rgb]{0.00,0.00,1.00}{##1}}}
\@namedef{PY@tok@vc}{\def\PY@tc##1{\textcolor[rgb]{0.10,0.09,0.49}{##1}}}
\@namedef{PY@tok@vg}{\def\PY@tc##1{\textcolor[rgb]{0.10,0.09,0.49}{##1}}}
\@namedef{PY@tok@vi}{\def\PY@tc##1{\textcolor[rgb]{0.10,0.09,0.49}{##1}}}
\@namedef{PY@tok@vm}{\def\PY@tc##1{\textcolor[rgb]{0.10,0.09,0.49}{##1}}}
\@namedef{PY@tok@sa}{\def\PY@tc##1{\textcolor[rgb]{0.73,0.13,0.13}{##1}}}
\@namedef{PY@tok@sb}{\def\PY@tc##1{\textcolor[rgb]{0.73,0.13,0.13}{##1}}}
\@namedef{PY@tok@sc}{\def\PY@tc##1{\textcolor[rgb]{0.73,0.13,0.13}{##1}}}
\@namedef{PY@tok@dl}{\def\PY@tc##1{\textcolor[rgb]{0.73,0.13,0.13}{##1}}}
\@namedef{PY@tok@s2}{\def\PY@tc##1{\textcolor[rgb]{0.73,0.13,0.13}{##1}}}
\@namedef{PY@tok@sh}{\def\PY@tc##1{\textcolor[rgb]{0.73,0.13,0.13}{##1}}}
\@namedef{PY@tok@s1}{\def\PY@tc##1{\textcolor[rgb]{0.73,0.13,0.13}{##1}}}
\@namedef{PY@tok@mb}{\def\PY@tc##1{\textcolor[rgb]{0.40,0.40,0.40}{##1}}}
\@namedef{PY@tok@mf}{\def\PY@tc##1{\textcolor[rgb]{0.40,0.40,0.40}{##1}}}
\@namedef{PY@tok@mh}{\def\PY@tc##1{\textcolor[rgb]{0.40,0.40,0.40}{##1}}}
\@namedef{PY@tok@mi}{\def\PY@tc##1{\textcolor[rgb]{0.40,0.40,0.40}{##1}}}
\@namedef{PY@tok@il}{\def\PY@tc##1{\textcolor[rgb]{0.40,0.40,0.40}{##1}}}
\@namedef{PY@tok@mo}{\def\PY@tc##1{\textcolor[rgb]{0.40,0.40,0.40}{##1}}}
\@namedef{PY@tok@ch}{\let\PY@it=\textit\def\PY@tc##1{\textcolor[rgb]{0.24,0.48,0.48}{##1}}}
\@namedef{PY@tok@cm}{\let\PY@it=\textit\def\PY@tc##1{\textcolor[rgb]{0.24,0.48,0.48}{##1}}}
\@namedef{PY@tok@cpf}{\let\PY@it=\textit\def\PY@tc##1{\textcolor[rgb]{0.24,0.48,0.48}{##1}}}
\@namedef{PY@tok@c1}{\let\PY@it=\textit\def\PY@tc##1{\textcolor[rgb]{0.24,0.48,0.48}{##1}}}
\@namedef{PY@tok@cs}{\let\PY@it=\textit\def\PY@tc##1{\textcolor[rgb]{0.24,0.48,0.48}{##1}}}

\def\PYZbs{\char`\\}
\def\PYZus{\char`\_}
\def\PYZob{\char`\{}
\def\PYZcb{\char`\}}
\def\PYZca{\char`\^}
\def\PYZam{\char`\&}
\def\PYZlt{\char`\<}
\def\PYZgt{\char`\>}
\def\PYZsh{\char`\#}
\def\PYZpc{\char`\%}
\def\PYZdl{\char`\$}
\def\PYZhy{\char`\-}
\def\PYZsq{\char`\'}
\def\PYZdq{\char`\"}
\def\PYZti{\char`\~}
% for compatibility with earlier versions
\def\PYZat{@}
\def\PYZlb{[}
\def\PYZrb{]}
\makeatother


    % For linebreaks inside Verbatim environment from package fancyvrb.
    \makeatletter
        \newbox\Wrappedcontinuationbox
        \newbox\Wrappedvisiblespacebox
        \newcommand*\Wrappedvisiblespace {\textcolor{red}{\textvisiblespace}}
        \newcommand*\Wrappedcontinuationsymbol {\textcolor{red}{\llap{\tiny$\m@th\hookrightarrow$}}}
        \newcommand*\Wrappedcontinuationindent {3ex }
        \newcommand*\Wrappedafterbreak {\kern\Wrappedcontinuationindent\copy\Wrappedcontinuationbox}
        % Take advantage of the already applied Pygments mark-up to insert
        % potential linebreaks for TeX processing.
        %        {, <, #, %, $, ' and ": go to next line.
        %        _, }, ^, &, >, - and ~: stay at end of broken line.
        % Use of \textquotesingle for straight quote.
        \newcommand*\Wrappedbreaksatspecials {%
            \def\PYGZus{\discretionary{\char`\_}{\Wrappedafterbreak}{\char`\_}}%
            \def\PYGZob{\discretionary{}{\Wrappedafterbreak\char`\{}{\char`\{}}%
            \def\PYGZcb{\discretionary{\char`\}}{\Wrappedafterbreak}{\char`\}}}%
            \def\PYGZca{\discretionary{\char`\^}{\Wrappedafterbreak}{\char`\^}}%
            \def\PYGZam{\discretionary{\char`\&}{\Wrappedafterbreak}{\char`\&}}%
            \def\PYGZlt{\discretionary{}{\Wrappedafterbreak\char`\<}{\char`\<}}%
            \def\PYGZgt{\discretionary{\char`\>}{\Wrappedafterbreak}{\char`\>}}%
            \def\PYGZsh{\discretionary{}{\Wrappedafterbreak\char`\#}{\char`\#}}%
            \def\PYGZpc{\discretionary{}{\Wrappedafterbreak\char`\%}{\char`\%}}%
            \def\PYGZdl{\discretionary{}{\Wrappedafterbreak\char`\$}{\char`\$}}%
            \def\PYGZhy{\discretionary{\char`\-}{\Wrappedafterbreak}{\char`\-}}%
            \def\PYGZsq{\discretionary{}{\Wrappedafterbreak\textquotesingle}{\textquotesingle}}%
            \def\PYGZdq{\discretionary{}{\Wrappedafterbreak\char`\"}{\char`\"}}%
            \def\PYGZti{\discretionary{\char`\~}{\Wrappedafterbreak}{\char`\~}}%
        }
        % Some characters . , ; ? ! / are not pygmentized.
        % This macro makes them "active" and they will insert potential linebreaks
        \newcommand*\Wrappedbreaksatpunct {%
            \lccode`\~`\.\lowercase{\def~}{\discretionary{\hbox{\char`\.}}{\Wrappedafterbreak}{\hbox{\char`\.}}}%
            \lccode`\~`\,\lowercase{\def~}{\discretionary{\hbox{\char`\,}}{\Wrappedafterbreak}{\hbox{\char`\,}}}%
            \lccode`\~`\;\lowercase{\def~}{\discretionary{\hbox{\char`\;}}{\Wrappedafterbreak}{\hbox{\char`\;}}}%
            \lccode`\~`\:\lowercase{\def~}{\discretionary{\hbox{\char`\:}}{\Wrappedafterbreak}{\hbox{\char`\:}}}%
            \lccode`\~`\?\lowercase{\def~}{\discretionary{\hbox{\char`\?}}{\Wrappedafterbreak}{\hbox{\char`\?}}}%
            \lccode`\~`\!\lowercase{\def~}{\discretionary{\hbox{\char`\!}}{\Wrappedafterbreak}{\hbox{\char`\!}}}%
            \lccode`\~`\/\lowercase{\def~}{\discretionary{\hbox{\char`\/}}{\Wrappedafterbreak}{\hbox{\char`\/}}}%
            \catcode`\.\active
            \catcode`\,\active
            \catcode`\;\active
            \catcode`\:\active
            \catcode`\?\active
            \catcode`\!\active
            \catcode`\/\active
            \lccode`\~`\~
        }
    \makeatother

    \let\OriginalVerbatim=\Verbatim
    \makeatletter
    \renewcommand{\Verbatim}[1][1]{%
        %\parskip\z@skip
        \sbox\Wrappedcontinuationbox {\Wrappedcontinuationsymbol}%
        \sbox\Wrappedvisiblespacebox {\FV@SetupFont\Wrappedvisiblespace}%
        \def\FancyVerbFormatLine ##1{\hsize\linewidth
            \vtop{\raggedright\hyphenpenalty\z@\exhyphenpenalty\z@
                \doublehyphendemerits\z@\finalhyphendemerits\z@
                \strut ##1\strut}%
        }%
        % If the linebreak is at a space, the latter will be displayed as visible
        % space at end of first line, and a continuation symbol starts next line.
        % Stretch/shrink are however usually zero for typewriter font.
        \def\FV@Space {%
            \nobreak\hskip\z@ plus\fontdimen3\font minus\fontdimen4\font
            \discretionary{\copy\Wrappedvisiblespacebox}{\Wrappedafterbreak}
            {\kern\fontdimen2\font}%
        }%

        % Allow breaks at special characters using \PYG... macros.
        \Wrappedbreaksatspecials
        % Breaks at punctuation characters . , ; ? ! and / need catcode=\active
        \OriginalVerbatim[#1,codes*=\Wrappedbreaksatpunct]%
    }
    \makeatother

    % Exact colors from NB
    \definecolor{incolor}{HTML}{303F9F}
    \definecolor{outcolor}{HTML}{D84315}
    \definecolor{cellborder}{HTML}{CFCFCF}
    \definecolor{cellbackground}{HTML}{F7F7F7}

    % prompt
    \makeatletter
    \newcommand{\boxspacing}{\kern\kvtcb@left@rule\kern\kvtcb@boxsep}
    \makeatother
    \newcommand{\prompt}[4]{
        {\ttfamily\llap{{\color{#2}[#3]:\hspace{3pt}#4}}\vspace{-\baselineskip}}
    }
    

    
    % Prevent overflowing lines due to hard-to-break entities
    \sloppy
    % Setup hyperref package
    \hypersetup{
      breaklinks=true,  % so long urls are correctly broken across lines
      colorlinks=true,
      urlcolor=urlcolor,
      linkcolor=linkcolor,
      citecolor=citecolor,
      }
    % Slightly bigger margins than the latex defaults
    
    \geometry{verbose,tmargin=1in,bmargin=1in,lmargin=1in,rmargin=1in}
    
    

\begin{document}
    
    \maketitle
    
    

    
    \textbf{Geração de Números Aleatórios em Python}

João Gabriel Miguêz da Silva - 2023008890 Felipe Bastos Vargas -
2023008881

\hypertarget{introduuxe7uxe3o}{%
\section{1.Introdução}\label{introduuxe7uxe3o}}

Este trabalho tem como objetivo testar de diferentes modelos de
distribuição para dados gerados aleatóriamente

Dentre as distribuições discretas serão aplicadas:

\begin{enumerate}
\def\labelenumi{\arabic{enumi}.}
\tightlist
\item
  \textbf{Uniforme Discreta}
\item
  \textbf{Bernoulli}
\item
  \textbf{Binomial}
\item
  \textbf{Poisson}
\end{enumerate}

Já as distribuições continuas serão:

\begin{enumerate}
\def\labelenumi{\arabic{enumi}.}
\tightlist
\item
  \textbf{Uniforme contínua}
\item
  \textbf{Normal}
\item
  \textbf{Exponencial}
\item
  \textbf{Gama}
\item
  \textbf{Beta}
\item
  \textbf{t-Student}
\end{enumerate}

    \hypertarget{metodologia}{%
\section{Metodologia}\label{metodologia}}

Inicialmente geramos aleatóriamente conjunto de dados usando a semente
aleatória 42 para cada distribuição.

    \begin{tcolorbox}[breakable, size=fbox, boxrule=1pt, pad at break*=1mm,colback=cellbackground, colframe=cellborder]
\prompt{In}{incolor}{2}{\boxspacing}
\begin{Verbatim}[commandchars=\\\{\}]
\PY{k+kn}{import}\PY{+w}{ }\PY{n+nn}{numpy}\PY{+w}{ }\PY{k}{as}\PY{+w}{ }\PY{n+nn}{np}
\PY{k+kn}{from}\PY{+w}{ }\PY{n+nn}{scipy} \PY{o}{.} \PY{n}{stats} \PY{k+kn}{import}\PY{+w}{ }\PY{n+nn}{t}
\PY{k+kn}{import}\PY{+w}{ }\PY{n+nn}{pandas}\PY{+w}{ }\PY{k}{as}\PY{+w}{ }\PY{n+nn}{pd}
\PY{k+kn}{import}\PY{+w}{ }\PY{n+nn}{math}

\PY{n}{np}\PY{o}{.}\PY{n}{random}\PY{o}{.}\PY{n}{seed} \PY{p}{(}\PY{l+m+mi}{42}\PY{p}{)}

\PY{c+c1}{\PYZsh{} Distribuicoes discretas}
\PY{n}{ddsUniformeDiscreta} \PY{o}{=} \PY{n}{np}\PY{o}{.}\PY{n}{random}\PY{o}{.}\PY{n}{randint}\PY{p}{(}\PY{l+m+mi}{0}\PY{p}{,}\PY{l+m+mi}{10}\PY{p}{,}\PY{n}{size}\PY{o}{=}\PY{l+m+mi}{1000}\PY{p}{)}         \PY{c+c1}{\PYZsh{}uniforme discreta}
\PY{n}{ddsBernoulli} \PY{o}{=} \PY{n}{np}\PY{o}{.}\PY{n}{random}\PY{o}{.}\PY{n}{binomial}\PY{p}{(}\PY{n}{n} \PY{o}{=}\PY{l+m+mi}{1}\PY{p}{,}\PY{n}{p} \PY{o}{=}\PY{l+m+mf}{0.5}\PY{p}{,}\PY{n}{size} \PY{o}{=}\PY{l+m+mi}{1000}\PY{p}{)}       \PY{c+c1}{\PYZsh{}Bernoulli}
\PY{n}{ddsBinomial} \PY{o}{=} \PY{n}{np}\PY{o}{.}\PY{n}{random}\PY{o}{.}\PY{n}{binomial}\PY{p}{(}\PY{n}{n} \PY{o}{=}\PY{l+m+mi}{10}\PY{p}{,}\PY{n}{p} \PY{o}{=}\PY{l+m+mf}{0.5}\PY{p}{,}\PY{n}{size} \PY{o}{=}\PY{l+m+mi}{1000}\PY{p}{)}       \PY{c+c1}{\PYZsh{}Binomial}
\PY{n}{ddsPoisson} \PY{o}{=} \PY{n}{np}\PY{o}{.}\PY{n}{random}\PY{o}{.}\PY{n}{poisson}\PY{p}{(}\PY{n}{lam}\PY{o}{=}\PY{l+m+mi}{5}\PY{p}{,}\PY{n}{size} \PY{o}{=}\PY{l+m+mi}{1000}\PY{p}{)}                \PY{c+c1}{\PYZsh{}Poisson}


\PY{c+c1}{\PYZsh{} Distribuicoes continuas}
\PY{n}{ddsUniformeContinua} \PY{o}{=} \PY{n}{np}\PY{o}{.}\PY{n}{random}\PY{o}{.}\PY{n}{uniform}\PY{p}{(}\PY{l+m+mi}{0}\PY{p}{,}\PY{l+m+mi}{1}\PY{p}{,}\PY{n}{size} \PY{o}{=}\PY{l+m+mi}{1000}\PY{p}{)}            \PY{c+c1}{\PYZsh{}uniforme continua}
\PY{n}{ddsNormal} \PY{o}{=} \PY{n}{np}\PY{o}{.}\PY{n}{random}\PY{o}{.}\PY{n}{normal}\PY{p}{(}\PY{l+m+mi}{0}\PY{p}{,}\PY{l+m+mi}{1}\PY{p}{,}\PY{n}{size} \PY{o}{=}\PY{l+m+mi}{1000}\PY{p}{)}                       \PY{c+c1}{\PYZsh{}normal}
\PY{n}{ddsExponencial} \PY{o}{=} \PY{n}{np}\PY{o}{.}\PY{n}{random}\PY{o}{.}\PY{n}{exponential}\PY{p}{(}\PY{n}{scale} \PY{o}{=}\PY{l+m+mi}{2}\PY{p}{,}\PY{n}{size} \PY{o}{=}\PY{l+m+mi}{1000}\PY{p}{)}         \PY{c+c1}{\PYZsh{}exponencial}
\PY{n}{ddsGamma} \PY{o}{=} \PY{n}{np}\PY{o}{.}\PY{n}{random}\PY{o}{.}\PY{n}{gamma}\PY{p}{(}\PY{n}{shape} \PY{o}{=}\PY{l+m+mi}{2}\PY{p}{,}\PY{n}{scale} \PY{o}{=}\PY{l+m+mi}{2}\PY{p}{,}\PY{n}{size} \PY{o}{=}\PY{l+m+mi}{1000}\PY{p}{)}           \PY{c+c1}{\PYZsh{}gama}
\PY{n}{ddsBeta} \PY{o}{=} \PY{n}{np}\PY{o}{.}\PY{n}{random}\PY{o}{.}\PY{n}{beta}\PY{p}{(} \PY{n}{a} \PY{o}{=}\PY{l+m+mi}{2}\PY{p}{,}\PY{n}{b} \PY{o}{=}\PY{l+m+mi}{5}\PY{p}{,}\PY{n}{size} \PY{o}{=}\PY{l+m+mi}{1000}\PY{p}{)}                        \PY{c+c1}{\PYZsh{}beta}
\PY{n}{ddsTStudent}\PY{o}{=} \PY{n}{t}\PY{o}{.}\PY{n}{rvs}\PY{p}{(}\PY{n}{df} \PY{o}{=}\PY{l+m+mi}{5}\PY{p}{,}\PY{n}{size} \PY{o}{=}\PY{l+m+mi}{1000}\PY{p}{)}                                   \PY{c+c1}{\PYZsh{} t\PYZhy{}Student}
 
\end{Verbatim}
\end{tcolorbox}

    Em seguida analisamos as frequências observadas dos dados de cada
conjunto. Para o conjunto de dados discretos, criamos um dicionário para
que fosse possível contabilizar a frequência através de um loop.

    \hypertarget{frequencias-observadas-dados-discretos}{%
\subsection{Frequencias observadas dados
discretos}\label{frequencias-observadas-dados-discretos}}

    \begin{tcolorbox}[breakable, size=fbox, boxrule=1pt, pad at break*=1mm,colback=cellbackground, colframe=cellborder]
\prompt{In}{incolor}{3}{\boxspacing}
\begin{Verbatim}[commandchars=\\\{\}]
\PY{c+c1}{\PYZsh{}Frequencias dos dados Discretos }
\PY{n}{frqObservadaUniformeDiscreta} \PY{o}{=} \PY{p}{\PYZob{}}\PY{p}{\PYZcb{}}
\PY{n}{frqObservadaBernoulli} \PY{o}{=} \PY{p}{\PYZob{}}\PY{p}{\PYZcb{}}
\PY{n}{frqObservadaBinomial} \PY{o}{=} \PY{p}{\PYZob{}}\PY{p}{\PYZcb{}}
\PY{n}{frqObservadaPoisson} \PY{o}{=} \PY{p}{\PYZob{}}\PY{p}{\PYZcb{}}

\PY{n}{conjuntoDeDadosDiscretos} \PY{o}{=} \PY{p}{[}\PY{p}{(}\PY{n}{ddsUniformeDiscreta}\PY{p}{,} \PY{n}{frqObservadaUniformeDiscreta}\PY{p}{)}\PY{p}{,} 
                            \PY{p}{(}\PY{n}{ddsBernoulli}\PY{p}{,} \PY{n}{frqObservadaBernoulli}\PY{p}{)}\PY{p}{,}
                              \PY{p}{(}\PY{n}{ddsBinomial}\PY{p}{,} \PY{n}{frqObservadaBinomial}\PY{p}{)}\PY{p}{,} 
                              \PY{p}{(}\PY{n}{ddsPoisson}\PY{p}{,}\PY{n}{frqObservadaPoisson}\PY{p}{)}\PY{p}{]}

\PY{k}{for} \PY{n}{dados}\PY{p}{,}\PY{n}{frq} \PY{o+ow}{in} \PY{n}{conjuntoDeDadosDiscretos}\PY{p}{:}
        \PY{k}{for} \PY{n}{valor} \PY{o+ow}{in} \PY{n}{dados}\PY{p}{:}
                \PY{k}{if} \PY{n}{valor} \PY{o+ow}{in} \PY{n}{frq}\PY{p}{:}
                        \PY{n}{frq}\PY{p}{[}\PY{n}{valor}\PY{p}{]} \PY{o}{+}\PY{o}{=} \PY{l+m+mi}{1}  
                \PY{k}{else}\PY{p}{:}
                        \PY{n}{frq}\PY{p}{[}\PY{n}{valor}\PY{p}{]} \PY{o}{=} \PY{l+m+mi}{1}
\end{Verbatim}
\end{tcolorbox}

    \hypertarget{frequencias-observadas-dados-continuos}{%
\subsection{Frequencias observadas dados
continuos}\label{frequencias-observadas-dados-continuos}}

    \begin{tcolorbox}[breakable, size=fbox, boxrule=1pt, pad at break*=1mm,colback=cellbackground, colframe=cellborder]
\prompt{In}{incolor}{4}{\boxspacing}
\begin{Verbatim}[commandchars=\\\{\}]
\PY{c+c1}{\PYZsh{}Continuas}
\PY{n}{frqObservadaUniformeContinua} \PY{o}{=} \PY{p}{\PYZob{}}\PY{p}{\PYZcb{}}
\PY{n}{frqObservadaNormal} \PY{o}{=} \PY{p}{\PYZob{}}\PY{p}{\PYZcb{}}
\PY{n}{frqObservadaExponencial} \PY{o}{=} \PY{p}{\PYZob{}}\PY{p}{\PYZcb{}}
\PY{n}{frqObservadaGamma} \PY{o}{=} \PY{p}{\PYZob{}}\PY{p}{\PYZcb{}}
\PY{n}{frqObservadaBeta} \PY{o}{=} \PY{p}{\PYZob{}}\PY{p}{\PYZcb{}}
\PY{n}{frqObservadaTStudent} \PY{o}{=} \PY{p}{\PYZob{}}\PY{p}{\PYZcb{}}

\PY{n}{conjuntoDeDadosContinuos} \PY{o}{=} \PY{p}{[}\PY{p}{(}\PY{n}{ddsUniformeContinua}\PY{p}{,}\PY{n}{frqObservadaUniformeContinua}\PY{p}{)}\PY{p}{,} 
                            \PY{p}{(}\PY{n}{ddsNormal}\PY{p}{,}\PY{n}{frqObservadaNormal}\PY{p}{)}\PY{p}{,} 
                            \PY{p}{(}\PY{n}{ddsExponencial}\PY{p}{,}\PY{n}{frqObservadaExponencial}\PY{p}{)}\PY{p}{,} 
                            \PY{p}{(}\PY{n}{ddsGamma}\PY{p}{,}\PY{n}{frqObservadaGamma}\PY{p}{)}\PY{p}{,} 
                            \PY{p}{(}\PY{n}{ddsBeta}\PY{p}{,}\PY{n}{frqObservadaBeta}\PY{p}{)}\PY{p}{,} 
                            \PY{p}{(}\PY{n}{ddsTStudent}\PY{p}{,}\PY{n}{frqObservadaTStudent}\PY{p}{)}\PY{p}{]}
\end{Verbatim}
\end{tcolorbox}

    \hypertarget{cuxe1lculo-das-frequuxeancias-esperadas-nas-distribuiuxe7uxf5es}{%
\subsection{Cálculo das frequências esperadas nas
distribuições}\label{cuxe1lculo-das-frequuxeancias-esperadas-nas-distribuiuxe7uxf5es}}

    \hypertarget{distribuiuxe7uxe3o-uniforme-discreta}{%
\subsubsection{Distribuição Uniforme
Discreta}\label{distribuiuxe7uxe3o-uniforme-discreta}}

    \begin{tcolorbox}[breakable, size=fbox, boxrule=1pt, pad at break*=1mm,colback=cellbackground, colframe=cellborder]
\prompt{In}{incolor}{5}{\boxspacing}
\begin{Verbatim}[commandchars=\\\{\}]
\PY{n}{n} \PY{o}{=} \PY{l+m+mi}{10}          \PY{c+c1}{\PYZsh{} número de valores possíveis}
\PY{n}{a} \PY{o}{=} \PY{l+m+mi}{0}           \PY{c+c1}{\PYZsh{} valor mínimo}
\PY{n}{b} \PY{o}{=} \PY{n}{a} \PY{o}{+} \PY{n}{n} \PY{o}{\PYZhy{}} \PY{l+m+mi}{1}   \PY{c+c1}{\PYZsh{} valor máximo possível (0 + 10 \PYZhy{} 1 = 9)}

\PY{k}{def}\PY{+w}{ }\PY{n+nf}{FMP\PYZus{}DiscretaUniforme}\PY{p}{(}\PY{n}{n}\PY{p}{,} \PY{n}{a}\PY{p}{,} \PY{n}{b}\PY{p}{)}\PY{p}{:}
    
    \PY{n}{valores\PYZus{}x} \PY{o}{=} \PY{p}{[}\PY{n}{i} \PY{k}{for} \PY{n}{i} \PY{o+ow}{in} \PY{n+nb}{range}\PY{p}{(}\PY{n}{a}\PY{p}{,} \PY{n}{b} \PY{o}{+} \PY{l+m+mi}{1}\PY{p}{)}\PY{p}{]}
    \PY{n}{probabilidade} \PY{o}{=} \PY{l+m+mi}{1} \PY{o}{/} \PY{n}{n}
    \PY{n}{probabilidades} \PY{o}{=} \PY{p}{[}\PY{n}{probabilidade} \PY{k}{for} \PY{n}{\PYZus{}} \PY{o+ow}{in} \PY{n+nb}{range}\PY{p}{(}\PY{n}{n}\PY{p}{)}\PY{p}{]}
    \PY{n}{fmp\PYZus{}resultado} \PY{o}{=} \PY{n+nb}{dict}\PY{p}{(}\PY{n+nb}{zip}\PY{p}{(}\PY{n}{valores\PYZus{}x}\PY{p}{,} \PY{n}{probabilidades}\PY{p}{)}\PY{p}{)}
    
    \PY{k}{return} \PY{n}{fmp\PYZus{}resultado}

\PY{k}{def}\PY{+w}{ }\PY{n+nf}{calcular\PYZus{}estatisticas\PYZus{}uniforme}\PY{p}{(}\PY{n}{a}\PY{p}{,} \PY{n}{b}\PY{p}{)}\PY{p}{:}
    \PY{n}{esperanca} \PY{o}{=} \PY{p}{(}\PY{n}{a} \PY{o}{+} \PY{n}{b}\PY{p}{)} \PY{o}{/} \PY{l+m+mi}{2}
    \PY{n}{variancia} \PY{o}{=} \PY{p}{(}\PY{p}{(}\PY{n}{b} \PY{o}{\PYZhy{}} \PY{n}{a} \PY{o}{+} \PY{l+m+mi}{1}\PY{p}{)}\PY{o}{*}\PY{o}{*}\PY{l+m+mi}{2} \PY{o}{\PYZhy{}} \PY{l+m+mi}{1}\PY{p}{)} \PY{o}{/} \PY{l+m+mi}{12}    
    \PY{k}{return} \PY{n}{esperanca}\PY{p}{,} \PY{n}{variancia}


\PY{n}{fmpDiscretaUniforme} \PY{o}{=} \PY{n}{FMP\PYZus{}DiscretaUniforme}\PY{p}{(}\PY{n}{n}\PY{p}{,} \PY{n}{a}\PY{p}{,} \PY{n}{b}\PY{p}{)}
\PY{n}{mediaTeoricaDiscUni}\PY{p}{,} \PY{n}{varTeoricaDiscUni} \PY{o}{=} \PY{n}{calcular\PYZus{}estatisticas\PYZus{}uniforme}\PY{p}{(}\PY{n}{a}\PY{p}{,} \PY{n}{b}\PY{p}{)}
\end{Verbatim}
\end{tcolorbox}

    \hypertarget{distribuiuxe7uxe3o-de-bernoulli}{%
\subsubsection{Distribuição de
Bernoulli}\label{distribuiuxe7uxe3o-de-bernoulli}}

    \begin{tcolorbox}[breakable, size=fbox, boxrule=1pt, pad at break*=1mm,colback=cellbackground, colframe=cellborder]
\prompt{In}{incolor}{6}{\boxspacing}
\begin{Verbatim}[commandchars=\\\{\}]
\PY{n}{pSucesso} \PY{o}{=} \PY{l+m+mf}{0.5}          
\PY{n}{pFracasso} \PY{o}{=} \PY{l+m+mi}{1} \PY{o}{\PYZhy{}} \PY{n}{pSucesso}

\PY{k}{def}\PY{+w}{ }\PY{n+nf}{FMP\PYZus{}Bernoulli}\PY{p}{(}\PY{n}{p}\PY{p}{)}\PY{p}{:}

    \PY{n}{valores\PYZus{}x} \PY{o}{=} \PY{p}{[}\PY{l+m+mi}{0}\PY{p}{,} \PY{l+m+mi}{1}\PY{p}{]}
    \PY{n}{probabilidades} \PY{o}{=} \PY{p}{[}\PY{l+m+mi}{1} \PY{o}{\PYZhy{}} \PY{n}{p}\PY{p}{,} \PY{n}{p}\PY{p}{]}
    \PY{n}{fmp\PYZus{}resultado} \PY{o}{=} \PY{n+nb}{dict}\PY{p}{(}\PY{n+nb}{zip}\PY{p}{(}\PY{n}{valores\PYZus{}x}\PY{p}{,} \PY{n}{probabilidades}\PY{p}{)}\PY{p}{)}
    \PY{k}{return} \PY{n}{fmp\PYZus{}resultado}

\PY{k}{def}\PY{+w}{ }\PY{n+nf}{calcular\PYZus{}estatisticas\PYZus{}bernoulli}\PY{p}{(}\PY{n}{p}\PY{p}{)}\PY{p}{:}
    \PY{n}{esperanca} \PY{o}{=} \PY{n}{p}
    \PY{n}{variancia} \PY{o}{=} \PY{n}{p} \PY{o}{*} \PY{p}{(}\PY{l+m+mi}{1} \PY{o}{\PYZhy{}} \PY{n}{p}\PY{p}{)}
    \PY{k}{return} \PY{n}{esperanca}\PY{p}{,} \PY{n}{variancia}

\PY{n}{fmpBernoulli} \PY{o}{=} \PY{n}{FMP\PYZus{}Bernoulli}\PY{p}{(}\PY{n}{pSucesso}\PY{p}{)}
\PY{n}{mediaTeoricaBernoulli}\PY{p}{,} \PY{n}{varTeoricaBernoulli} \PY{o}{=} \PY{n}{calcular\PYZus{}estatisticas\PYZus{}bernoulli}\PY{p}{(}\PY{n}{pSucesso}\PY{p}{)}
\end{Verbatim}
\end{tcolorbox}

    \hypertarget{distribuiuxe7uxe3o-binomial}{%
\subsubsection{Distribuição
Binomial}\label{distribuiuxe7uxe3o-binomial}}

    \begin{tcolorbox}[breakable, size=fbox, boxrule=1pt, pad at break*=1mm,colback=cellbackground, colframe=cellborder]
\prompt{In}{incolor}{7}{\boxspacing}
\begin{Verbatim}[commandchars=\\\{\}]
\PY{n}{n} \PY{o}{=} \PY{l+m+mi}{10}   \PY{c+c1}{\PYZsh{} Número de tentativas independentes}
\PY{n}{p} \PY{o}{=} \PY{l+m+mf}{0.5}  \PY{c+c1}{\PYZsh{} Probabilidade de sucesso em cada tentativa}

\PY{k}{def}\PY{+w}{ }\PY{n+nf}{FMP\PYZus{}Binomial}\PY{p}{(}\PY{n}{n}\PY{p}{,} \PY{n}{p}\PY{p}{)}\PY{p}{:}
    \PY{n}{fmp\PYZus{}resultado} \PY{o}{=} \PY{p}{\PYZob{}}\PY{p}{\PYZcb{}}
    \PY{k}{for} \PY{n}{k} \PY{o+ow}{in} \PY{n+nb}{range}\PY{p}{(}\PY{n}{n} \PY{o}{+} \PY{l+m+mi}{1}\PY{p}{)}\PY{p}{:}
        \PY{n}{coeficiente} \PY{o}{=} \PY{n}{math}\PY{o}{.}\PY{n}{comb}\PY{p}{(}\PY{n}{n}\PY{p}{,} \PY{n}{k}\PY{p}{)}                               \PY{c+c1}{\PYZsh{} coeficiente binomial: n! / (k! * (n \PYZhy{} k)!)}
        \PY{n}{probabilidade} \PY{o}{=} \PY{n}{coeficiente} \PY{o}{*} \PY{p}{(}\PY{n}{p}\PY{o}{*}\PY{o}{*}\PY{n}{k}\PY{p}{)} \PY{o}{*} \PY{p}{(}\PY{p}{(}\PY{l+m+mi}{1} \PY{o}{\PYZhy{}} \PY{n}{p}\PY{p}{)}\PY{o}{*}\PY{o}{*}\PY{p}{(}\PY{n}{n} \PY{o}{\PYZhy{}} \PY{n}{k}\PY{p}{)}\PY{p}{)}   \PY{c+c1}{\PYZsh{}P(X = k) = Coeficiente * p\PYZca{}k * (1\PYZhy{}p)\PYZca{}(n\PYZhy{}k)}
        \PY{c+c1}{\PYZsh{} Guardar no dicionário}
        \PY{n}{fmp\PYZus{}resultado}\PY{p}{[}\PY{n}{k}\PY{p}{]} \PY{o}{=} \PY{n}{probabilidade}
    \PY{k}{return} \PY{n}{fmp\PYZus{}resultado}


\PY{k}{def}\PY{+w}{ }\PY{n+nf}{calcular\PYZus{}estatisticas\PYZus{}binomial}\PY{p}{(}\PY{n}{n}\PY{p}{,} \PY{n}{p}\PY{p}{)}\PY{p}{:}
    \PY{n}{esperanca} \PY{o}{=} \PY{n}{n} \PY{o}{*} \PY{n}{p}
    \PY{n}{variancia} \PY{o}{=} \PY{n}{n} \PY{o}{*} \PY{n}{p} \PY{o}{*} \PY{p}{(}\PY{l+m+mi}{1} \PY{o}{\PYZhy{}} \PY{n}{p}\PY{p}{)}
    \PY{k}{return} \PY{n}{esperanca}\PY{p}{,} \PY{n}{variancia}

\PY{n}{fmpBinomial} \PY{o}{=} \PY{n}{FMP\PYZus{}Binomial}\PY{p}{(}\PY{n}{n}\PY{p}{,} \PY{n}{p}\PY{p}{)}
\PY{n}{mediaTeoricaBinomial}\PY{p}{,} \PY{n}{varTeoricaBinomial} \PY{o}{=} \PY{n}{calcular\PYZus{}estatisticas\PYZus{}binomial}\PY{p}{(}\PY{n}{n}\PY{p}{,} \PY{n}{p}\PY{p}{)}
\end{Verbatim}
\end{tcolorbox}

    \hypertarget{distribuiuxe7uxe3o-poisson}{%
\subsubsection{Distribuição Poisson}\label{distribuiuxe7uxe3o-poisson}}

    \begin{tcolorbox}[breakable, size=fbox, boxrule=1pt, pad at break*=1mm,colback=cellbackground, colframe=cellborder]
\prompt{In}{incolor}{8}{\boxspacing}
\begin{Verbatim}[commandchars=\\\{\}]
\PY{n}{lam} \PY{o}{=} \PY{l+m+mi}{5}  \PY{c+c1}{\PYZsh{} Lambda}

\PY{k}{def}\PY{+w}{ }\PY{n+nf}{FMP\PYZus{}Poisson}\PY{p}{(}\PY{n}{lam}\PY{p}{,} \PY{n}{limite\PYZus{}max}\PY{o}{=}\PY{l+m+mi}{15}\PY{p}{)}\PY{p}{:}
    \PY{n}{fmp\PYZus{}resultado} \PY{o}{=} \PY{p}{\PYZob{}}\PY{p}{\PYZcb{}}
    \PY{k}{for} \PY{n}{k} \PY{o+ow}{in} \PY{n+nb}{range}\PY{p}{(}\PY{n}{limite\PYZus{}max} \PY{o}{+} \PY{l+m+mi}{1}\PY{p}{)}\PY{p}{:}

        \PY{n}{fatorial\PYZus{}k} \PY{o}{=} \PY{n}{math}\PY{o}{.}\PY{n}{factorial}\PY{p}{(}\PY{n}{k}\PY{p}{)}
        \PY{n}{probabilidade} \PY{o}{=} \PY{p}{(}\PY{n}{math}\PY{o}{.}\PY{n}{exp}\PY{p}{(}\PY{o}{\PYZhy{}}\PY{n}{lam}\PY{p}{)} \PY{o}{*} \PY{p}{(}\PY{n}{lam}\PY{o}{*}\PY{o}{*}\PY{n}{k}\PY{p}{)}\PY{p}{)} \PY{o}{/} \PY{n}{fatorial\PYZus{}k}    \PY{c+c1}{\PYZsh{}P(X = k) = (e\PYZca{}(\PYZhy{}λ) * λ\PYZca{}k) / k!}
    
        \PY{c+c1}{\PYZsh{} Guardar no dicionário}
        \PY{n}{fmp\PYZus{}resultado}\PY{p}{[}\PY{n}{k}\PY{p}{]} \PY{o}{=} \PY{n}{probabilidade}
    \PY{k}{return} \PY{n}{fmp\PYZus{}resultado}

\PY{k}{def}\PY{+w}{ }\PY{n+nf}{calcular\PYZus{}estatisticas\PYZus{}poisson}\PY{p}{(}\PY{n}{lam}\PY{p}{)}\PY{p}{:}
    \PY{k}{return}  \PY{n}{lam}\PY{p}{,} \PY{n}{lam}

\PY{n}{fpmPoisson} \PY{o}{=} \PY{n}{FMP\PYZus{}Poisson}\PY{p}{(}\PY{n}{lam}\PY{p}{)}
\PY{n}{mediaTeoriaPoisson}\PY{p}{,} \PY{n}{varTeoricaPoisson} \PY{o}{=} \PY{n}{calcular\PYZus{}estatisticas\PYZus{}poisson}\PY{p}{(}\PY{n}{lam}\PY{p}{)}
\end{Verbatim}
\end{tcolorbox}

    \hypertarget{analise-dos-resultados}{%
\section{Analise dos Resultados}\label{analise-dos-resultados}}

    \begin{tcolorbox}[breakable, size=fbox, boxrule=1pt, pad at break*=1mm,colback=cellbackground, colframe=cellborder]
\prompt{In}{incolor}{ }{\boxspacing}
\begin{Verbatim}[commandchars=\\\{\}]
\PY{n}{tabelas\PYZus{}frequencia} \PY{o}{=} \PY{p}{[}
    \PY{p}{(}\PY{l+s+s2}{\PYZdq{}}\PY{l+s+s2}{Uniforme Discreta}\PY{l+s+s2}{\PYZdq{}}\PY{p}{,} \PY{n}{frqObservadaUniformeDiscreta}\PY{p}{,} \PY{n}{fmpDiscretaUniforme}\PY{p}{)}\PY{p}{,}
    \PY{p}{(}\PY{l+s+s2}{\PYZdq{}}\PY{l+s+s2}{Bernoulli}\PY{l+s+s2}{\PYZdq{}}\PY{p}{,} \PY{n}{frqObservadaBernoulli}\PY{p}{,} \PY{n}{fmpBernoulli}\PY{p}{)}\PY{p}{,}
    \PY{p}{(}\PY{l+s+s2}{\PYZdq{}}\PY{l+s+s2}{Binomial}\PY{l+s+s2}{\PYZdq{}}\PY{p}{,} \PY{n}{frqObservadaBinomial}\PY{p}{,} \PY{n}{fmpBinomial}\PY{p}{)}\PY{p}{,}
    \PY{p}{(}\PY{l+s+s2}{\PYZdq{}}\PY{l+s+s2}{Poisson}\PY{l+s+s2}{\PYZdq{}}\PY{p}{,} \PY{n}{frqObservadaPoisson}\PY{p}{,} \PY{n}{fpmPoisson}\PY{p}{)}
\PY{p}{]}

\PY{k}{for} \PY{n}{nome}\PY{p}{,} \PY{n}{frq\PYZus{}dict}\PY{p}{,} \PY{n}{fmp\PYZus{}dict} \PY{o+ow}{in} \PY{n}{tabelas\PYZus{}frequencia}\PY{p}{:}
    \PY{n+nb}{print}\PY{p}{(}\PY{l+s+sa}{f}\PY{l+s+s2}{\PYZdq{}}\PY{l+s+se}{\PYZbs{}n}\PY{l+s+s2}{\PYZdq{}} \PY{o}{+} \PY{l+s+s2}{\PYZdq{}}\PY{l+s+s2}{=}\PY{l+s+s2}{\PYZdq{}}\PY{o}{*}\PY{l+m+mi}{70}\PY{p}{)}
    \PY{n+nb}{print}\PY{p}{(}\PY{l+s+sa}{f}\PY{l+s+s2}{\PYZdq{}}\PY{l+s+s2}{ TABELA COMPARATIVA: }\PY{l+s+si}{\PYZob{}}\PY{n}{nome}\PY{o}{.}\PY{n}{upper}\PY{p}{(}\PY{p}{)}\PY{l+s+si}{\PYZcb{}}\PY{l+s+s2}{ }\PY{l+s+s2}{\PYZdq{}}\PY{p}{)}
    \PY{n+nb}{print}\PY{p}{(}\PY{l+s+s2}{\PYZdq{}}\PY{l+s+s2}{=}\PY{l+s+s2}{\PYZdq{}}\PY{o}{*}\PY{l+m+mi}{70}\PY{p}{)}
    
    
    \PY{n}{df} \PY{o}{=} \PY{n}{pd}\PY{o}{.}\PY{n}{DataFrame}\PY{p}{(}\PY{n+nb}{list}\PY{p}{(}\PY{n}{frq\PYZus{}dict}\PY{o}{.}\PY{n}{items}\PY{p}{(}\PY{p}{)}\PY{p}{)}\PY{p}{,} \PY{n}{columns}\PY{o}{=}\PY{p}{[}\PY{l+s+s1}{\PYZsq{}}\PY{l+s+s1}{Valor}\PY{l+s+s1}{\PYZsq{}}\PY{p}{,} \PY{l+s+s1}{\PYZsq{}}\PY{l+s+s1}{Freq. Obs. Absoluta (f\PYZus{}o)}\PY{l+s+s1}{\PYZsq{}}\PY{p}{]}\PY{p}{)}
    \PY{n}{df} \PY{o}{=} \PY{n}{df}\PY{o}{.}\PY{n}{sort\PYZus{}values}\PY{p}{(}\PY{n}{by}\PY{o}{=}\PY{l+s+s1}{\PYZsq{}}\PY{l+s+s1}{Valor}\PY{l+s+s1}{\PYZsq{}}\PY{p}{)}\PY{o}{.}\PY{n}{reset\PYZus{}index}\PY{p}{(}\PY{n}{drop}\PY{o}{=}\PY{k+kc}{True}\PY{p}{)}
    
    \PY{n}{total\PYZus{}observacoes} \PY{o}{=} \PY{n}{df}\PY{p}{[}\PY{l+s+s1}{\PYZsq{}}\PY{l+s+s1}{Freq. Obs. Absoluta (f\PYZus{}o)}\PY{l+s+s1}{\PYZsq{}}\PY{p}{]}\PY{o}{.}\PY{n}{sum}\PY{p}{(}\PY{p}{)}
    
    \PY{n}{df}\PY{p}{[}\PY{l+s+s1}{\PYZsq{}}\PY{l+s+s1}{Prob. Teórica P(X=x)}\PY{l+s+s1}{\PYZsq{}}\PY{p}{]} \PY{o}{=} \PY{n}{df}\PY{p}{[}\PY{l+s+s1}{\PYZsq{}}\PY{l+s+s1}{Valor}\PY{l+s+s1}{\PYZsq{}}\PY{p}{]}\PY{o}{.}\PY{n}{map}\PY{p}{(}\PY{n}{fmp\PYZus{}dict}\PY{p}{)}
    \PY{n}{df}\PY{p}{[}\PY{l+s+s1}{\PYZsq{}}\PY{l+s+s1}{Prob. Teórica P(X=x)}\PY{l+s+s1}{\PYZsq{}}\PY{p}{]} \PY{o}{=} \PY{n}{df}\PY{p}{[}\PY{l+s+s1}{\PYZsq{}}\PY{l+s+s1}{Prob. Teórica P(X=x)}\PY{l+s+s1}{\PYZsq{}}\PY{p}{]}\PY{o}{.}\PY{n}{fillna}\PY{p}{(}\PY{l+m+mi}{0}\PY{p}{)}
    \PY{n}{df}\PY{p}{[}\PY{l+s+s1}{\PYZsq{}}\PY{l+s+s1}{Freq. Esp. Absoluta (f\PYZus{}e)}\PY{l+s+s1}{\PYZsq{}}\PY{p}{]} \PY{o}{=} \PY{n}{df}\PY{p}{[}\PY{l+s+s1}{\PYZsq{}}\PY{l+s+s1}{Prob. Teórica P(X=x)}\PY{l+s+s1}{\PYZsq{}}\PY{p}{]} \PY{o}{*} \PY{n}{total\PYZus{}observacoes}
    \PY{n}{df}\PY{p}{[}\PY{l+s+s1}{\PYZsq{}}\PY{l+s+s1}{Freq. Obs. Relativa (f\PYZus{}r)}\PY{l+s+s1}{\PYZsq{}}\PY{p}{]} \PY{o}{=} \PY{n}{df}\PY{p}{[}\PY{l+s+s1}{\PYZsq{}}\PY{l+s+s1}{Freq. Obs. Absoluta (f\PYZus{}o)}\PY{l+s+s1}{\PYZsq{}}\PY{p}{]} \PY{o}{/} \PY{n}{total\PYZus{}observacoes}
    

    \PY{n}{df\PYZus{}formatado} \PY{o}{=} \PY{n}{df}\PY{o}{.}\PY{n}{copy}\PY{p}{(}\PY{p}{)}
    
    \PY{n}{df\PYZus{}formatado}\PY{p}{[}\PY{l+s+s1}{\PYZsq{}}\PY{l+s+s1}{Freq. Esp. Absoluta (f\PYZus{}e)}\PY{l+s+s1}{\PYZsq{}}\PY{p}{]} \PY{o}{=} \PY{n}{df\PYZus{}formatado}\PY{p}{[}\PY{l+s+s1}{\PYZsq{}}\PY{l+s+s1}{Freq. Esp. Absoluta (f\PYZus{}e)}\PY{l+s+s1}{\PYZsq{}}\PY{p}{]}\PY{o}{.}\PY{n}{map}\PY{p}{(}\PY{l+s+s2}{\PYZdq{}}\PY{l+s+si}{\PYZob{}:.1f\PYZcb{}}\PY{l+s+s2}{\PYZdq{}}\PY{o}{.}\PY{n}{format}\PY{p}{)}
    \PY{n}{df\PYZus{}formatado}\PY{p}{[}\PY{l+s+s1}{\PYZsq{}}\PY{l+s+s1}{Prob. Teórica P(X=x)}\PY{l+s+s1}{\PYZsq{}}\PY{p}{]} \PY{o}{=} \PY{n}{df\PYZus{}formatado}\PY{p}{[}\PY{l+s+s1}{\PYZsq{}}\PY{l+s+s1}{Prob. Teórica P(X=x)}\PY{l+s+s1}{\PYZsq{}}\PY{p}{]}\PY{o}{.}\PY{n}{map}\PY{p}{(}\PY{l+s+s2}{\PYZdq{}}\PY{l+s+si}{\PYZob{}:.2\PYZpc{}\PYZcb{}}\PY{l+s+s2}{\PYZdq{}}\PY{o}{.}\PY{n}{format}\PY{p}{)}
    \PY{n}{df\PYZus{}formatado}\PY{p}{[}\PY{l+s+s1}{\PYZsq{}}\PY{l+s+s1}{Freq. Obs. Relativa (f\PYZus{}r)}\PY{l+s+s1}{\PYZsq{}}\PY{p}{]} \PY{o}{=} \PY{n}{df\PYZus{}formatado}\PY{p}{[}\PY{l+s+s1}{\PYZsq{}}\PY{l+s+s1}{Freq. Obs. Relativa (f\PYZus{}r)}\PY{l+s+s1}{\PYZsq{}}\PY{p}{]}\PY{o}{.}\PY{n}{map}\PY{p}{(}\PY{l+s+s2}{\PYZdq{}}\PY{l+s+si}{\PYZob{}:.2\PYZpc{}\PYZcb{}}\PY{l+s+s2}{\PYZdq{}}\PY{o}{.}\PY{n}{format}\PY{p}{)}
    
    \PY{n}{colunas\PYZus{}ordenadas} \PY{o}{=} \PY{p}{[}
        \PY{l+s+s1}{\PYZsq{}}\PY{l+s+s1}{Valor}\PY{l+s+s1}{\PYZsq{}}\PY{p}{,} 
        \PY{l+s+s1}{\PYZsq{}}\PY{l+s+s1}{Freq. Obs. Absoluta (f\PYZus{}o)}\PY{l+s+s1}{\PYZsq{}}\PY{p}{,} \PY{l+s+s1}{\PYZsq{}}\PY{l+s+s1}{Freq. Esp. Absoluta (f\PYZus{}e)}\PY{l+s+s1}{\PYZsq{}}\PY{p}{,} 
        \PY{l+s+s1}{\PYZsq{}}\PY{l+s+s1}{Freq. Obs. Relativa (f\PYZus{}r)}\PY{l+s+s1}{\PYZsq{}}\PY{p}{,} \PY{l+s+s1}{\PYZsq{}}\PY{l+s+s1}{Prob. Teórica P(X=x)}\PY{l+s+s1}{\PYZsq{}}
    \PY{p}{]}
    \PY{n+nb}{print}\PY{p}{(}\PY{n}{df\PYZus{}formatado}\PY{p}{[}\PY{n}{colunas\PYZus{}ordenadas}\PY{p}{]}\PY{o}{.}\PY{n}{to\PYZus{}string}\PY{p}{(}\PY{n}{index}\PY{o}{=}\PY{k+kc}{False}\PY{p}{)}\PY{p}{)}
\end{Verbatim}
\end{tcolorbox}

    \begin{Verbatim}[commandchars=\\\{\}]

======================================================================
 TABELA COMPARATIVA: UNIFORME DISCRETA
======================================================================
 Valor  Freq. Obs. Absoluta (f\_o) Freq. Esp. Absoluta (f\_e) Freq. Obs. Relativa
(f\_r) Prob. Teórica P(X=x)
     0                        118                     100.0
11.80\%               10.00\%
     1                         83                     100.0
8.30\%               10.00\%
     2                        110                     100.0
11.00\%               10.00\%
     3                         94                     100.0
9.40\%               10.00\%
     4                        107                     100.0
10.70\%               10.00\%
     5                         96                     100.0
9.60\%               10.00\%
     6                         94                     100.0
9.40\%               10.00\%
     7                        100                     100.0
10.00\%               10.00\%
     8                         91                     100.0
9.10\%               10.00\%
     9                        107                     100.0
10.70\%               10.00\%

======================================================================
 TABELA COMPARATIVA: BERNOULLI
======================================================================
 Valor  Freq. Obs. Absoluta (f\_o) Freq. Esp. Absoluta (f\_e) Freq. Obs. Relativa
(f\_r) Prob. Teórica P(X=x)
     0                        489                     500.0
48.90\%               50.00\%
     1                        511                     500.0
51.10\%               50.00\%

======================================================================
 TABELA COMPARATIVA: BINOMIAL
======================================================================
 Valor  Freq. Obs. Absoluta (f\_o) Freq. Esp. Absoluta (f\_e) Freq. Obs. Relativa
(f\_r) Prob. Teórica P(X=x)
     0                          1                       1.0
0.10\%                0.10\%
     1                         12                       9.8
1.20\%                0.98\%
     2                         50                      43.9
5.00\%                4.39\%
     3                        125                     117.2
12.50\%               11.72\%
     4                        175                     205.1
17.50\%               20.51\%
     5                        246                     246.1
24.60\%               24.61\%
     6                        218                     205.1
21.80\%               20.51\%
     7                        115                     117.2
11.50\%               11.72\%
     8                         44                      43.9
4.40\%                4.39\%
     9                         12                       9.8
1.20\%                0.98\%
    10                          2                       1.0
0.20\%                0.10\%

======================================================================
 TABELA COMPARATIVA: POISSON
======================================================================
 Valor  Freq. Obs. Absoluta (f\_o) Freq. Esp. Absoluta (f\_e) Freq. Obs. Relativa
(f\_r) Prob. Teórica P(X=x)
     0                          7                       6.7
0.70\%                0.67\%
     1                         31                      33.7
3.10\%                3.37\%
     2                         88                      84.2
8.80\%                8.42\%
     3                        147                     140.4
14.70\%               14.04\%
     4                        177                     175.5
17.70\%               17.55\%
     5                        173                     175.5
17.30\%               17.55\%
     6                        151                     146.2
15.10\%               14.62\%
     7                         84                     104.4
8.40\%               10.44\%
     8                         67                      65.3
6.70\%                6.53\%
     9                         45                      36.3
4.50\%                3.63\%
    10                         15                      18.1
1.50\%                1.81\%
    11                          8                       8.2
0.80\%                0.82\%
    12                          5                       3.4
0.50\%                0.34\%
    13                          1                       1.3
0.10\%                0.13\%
    15                          1                       0.2
0.10\%                0.02\%
    \end{Verbatim}


    % Add a bibliography block to the postdoc
    
    
    
\end{document}
